\begin{frame}{Literature Review: Overview}
\begin{block}{Focus of the review}
This review examines how multi-agent LLM systems communicate, share knowledge and maintain useful memory across tasks and sessions.
\end{block}
\begin{columns}[T,onlytextwidth]
\begin{column}{0.48\textwidth}
\begin{block}{Themes across the literature}
\begin{itemize}
\item \textbf{Memory formation:} deciding what experience is worth retaining.
\item \textbf{Shared memory:} enabling agents to access common knowledge safely.
\item \textbf{Architecture:} organizing memory hierarchy, access and consistency.
\item \textbf{Governance:} controlling scope, provenance, ownership and updates.
\end{itemize}
\end{block}
\end{column}
\begin{column}{0.48\textwidth}
\begin{block}{Six-paper review path}
\begin{itemize}
\item Selective persistent memory
\item Governed shared memory
\item Computer-architecture perspective
\item Mesh memory protocol
\item Adaptive memory lifecycle
\item Change-intent and write governance
\end{itemize}
\end{block}
\end{column}
\end{columns}
\vspace{2mm}
\begin{alertblock}{Review direction}
The detailed review moves from communication and shared state to selective admission, provenance, conflict handling and long-term memory evolution. The final summary compares the papers and identifies the gap addressed by SYNAPSE-MAS.
\end{alertblock}
\end{frame}
